Optimization is such an important tool that it deserves a set of notes in itself. All problems in model training essentially stems from non-ideal optimization. Knowing the strengths and weaknesses of each optimizer allows you to diagnose which ones to use. 

Generally, we (non-exclusively) categorize optimization algorithms as such: 
\begin{enumerate}
  \item \textit{Convex}? Convex optimization is pretty easy to solve and has been studied extensively. For nonconvex optimization, none of the algorithms can guarantee that we will find the global minima, and this is one of the hardest problems in statistics.\footnote{In practice, when we are doing high-dimensional nonconvex optimization, the best we can do is play around with some properties. In these cases, 0th order approximations are hopeless since the dimensions are too high, and second order approximations are hopeless either since computing the Hessian is too expensive for one run. Therefore, we must resort to some first order methods. } 
  \item \textit{Constrained}? Is the parameter space constrained to a certain manifold? 
  \item \textit{Order}. Do we use derivatives at all? First-order derivatives (gradient)? Second-order (Hessian)? 
\end{enumerate}

These algorithms try to solve the following potential problems. 
\begin{enumerate}
  \item \textit{Convergence}. Do we converge to some point? 
  \item \textit{Optimality}. Is this point close to the true global minima?  
  \item \textit{Efficiency}. Can we iterate efficiently? 
\end{enumerate} 

In optimization, we are given a function $f: S \subset \mathbb{R}^d \to \mathbb{R}$, and we want to find 
\begin{equation}
	\inf_{w \in S} f
\end{equation}
Ideally, $S$ and $f$ would have some nice properties such that the minimizer actually exists, and we can actually find $w^\ast = \argmin_{w \in S} f$ but this may not be required all the time. 

\begin{definition}[Function Class]
	Given some domain $S \subset \mathbb{R}$, a \textbf{function class} is some set of functions 
	\begin{equation}
		\mathcal{F} \subset \{ f: S \to \mathbb{R} \}
	\end{equation}
\end{definition}

Given a function $f \in \mathcal{F}$, we would like to formalize extracting some information out of this function, such as its value, derivative, or whether it is continuous at a point. We can describe this in terms of an oracle. 

\begin{definition}[Oracle]
	Let $\mathcal{Y}$ be an arbitrary space.\footnote{It is a space where the information that we wish to query lives.} An \textbf{oracle} 
	\begin{enumerate}
		\item for a function $f: S \to \mathbb{R}$ is a map
			\begin{equation}
				\mathcal{O}_f: S \to \mathcal{Y}
			\end{equation}

		\item for a function class $\mathcal{F} = \{f : S \to \mathbb{R}\}$ is a family of oracles 
			\begin{equation}
				\mathcal{O} = \{\mathcal{O}_f\}_{f \in \mathcal{F}}, \quad \mathcal{O}_f: S \to \mathcal{Y}
			\end{equation}
	\end{enumerate}
\end{definition}

Note that an oracle can be defined for a particular function or a function class. 

\begin{example}[First Order Oracle]
	If $\mathcal{F}$ be the space of functions that are differentiable in an open neighborhood of $S$. The first-order oracle is 
	\begin{equation}
		\mathcal{O}_f (x) = \big( f(x), \nabla f(x) \big)
	\end{equation}
	Therefore, it allows us to query both the value of the function and the gradient at that point. In practice, we do this using \textit{automatic differentiation}. 
\end{example}

\begin{definition}[Algorithm]
	Let $\mathcal{F}$ be a function class over $S \subset \mathbb{R}^d$ and fix $x_0 \in S$. A \textbf{deterministic $k$-query oracle algorithm} for $f \in \mathcal{F}$ is an ordered tuple $\mathcal{A}_k = (\Phi_1, \ldots, \Phi_{k-1}, \Psi_k)$ consisting of 
	\begin{enumerate}
	  \item \textit{Query Rules}. A function that decides where the algorithm should evaluate the oracle next. 
			\begin{equation}
				\Phi_t : (S \times \mathcal{Y})^t \to S, \quad t = 1, \ldots, k-1
			\end{equation}
			defined recursively as 
			\begin{align}
				y_t & = \mathcal{O}_f (x_t) \\ 
				x_{t+1} & = \Phi_{t+1} (x_0, y_0, \ldots, x_t, y_t)
			\end{align}

		\item \textit{Output Rule}. A function that decides which point to return after collecting $k$ oracle responses.\footnote{We can think of $k$ as the number of iterations.} 
			\begin{equation}
				\Psi_k : (S \times \mathcal{Y})^k \to S
			\end{equation}
			and the \textbf{output} of the algorithm is defined to be 
			\begin{equation}
				w_{\mathcal{A}_k} (f) = \Psi_k (x_0, y_0, \ldots, x_{k-1}, y_{k-1})
			\end{equation}
	\end{enumerate}
\end{definition}

\begin{example}[Gradient Descent]
	Let $\mathcal{F}$ be the class of differentiable functions $f: \mathbb{R}^d \to \mathbb{R}$. We can use the fact that \hyperref[mul-thm:gradient_increase]{gradients point in the step of steepest ascent}. This inspires a greedy approach to simply walk in this direction $-\nabla f(x)$ to minimize a function, so to formalize this, consider the first-order oracle 
	\begin{equation}
		\mathcal{O}_f : S \to \mathbb{R} \times \mathbb{R}^d, \quad \mathcal{O}_f (x) = \big( f(x), \nabla f(x) \big)
	\end{equation}
	Fix an initial point $x_0 \in S$ and a sequence of positive step sizes $(\eta_t)$. \textbf{Gradient descent} is the first-order deterministic algorithm defined by the update maps 
	\begin{equation}
		\Phi_t : (S \times (\mathbb{R} \times \mathbb{R}^d))^{t+1} \to S, \quad \Phi_t (x_0, y_0, \ldots, x_t, y_t) = x_t - \eta_t \pi_2 (y_t) = x_t - \eta_t \nabla f(x_t)
	\end{equation}
	where $\pi_2 : \mathbb{R} \times \mathbb{R}^d \to \mathbb{R}^d$ is projection onto the gradient component. Therefore, 
	\begin{equation}
		x_{t+1} = x_t - \eta_t \nabla f(x_t), \quad t = 0, \ldots, k-1
	\end{equation}
	After $k$ iterations, the algorithm output $w_{\mathcal{A}_k} (f) = x_k$. 
\end{example}

\begin{definition}[Performance Criterion]
	Given a function class $\mathcal{F}$ defined over $S$, a \textbf{performance criterion} is a function 
	\begin{equation}
		\mathcal{L}: \mathcal{F} \times S \to [0, +\infty]
	\end{equation}
	Sometimes, we write $\mathcal{L}_f (w) = \mathcal{L}(f, w)$. For minimizations, the performance criterion should ideally satisfy $\mathcal{L}_f (w^\ast) = 0$ for some optimal value $w^\ast$. 
\end{definition}

There are many metrics that we can use to measure optimality, in increasing importance. 

\begin{enumerate}
	\item \textit{Small Gradient Norm}. How small is $\| \nabla f(w)\|$? By first-order optimality conditions, a local minimum must have vanishing gradient, so perhaps bounding the norm gradient gives us a few clues. 

	\item \textit{Small Objective Gap}. Maybe we wish to bound 
		\begin{equation}
			f(w) - f(w^\ast) 
		\end{equation}

	\item \textit{Small Distance to Solution}. Maybe we wish to bound 
		\begin{equation}
			\|w - w^\ast\|
		\end{equation}
\end{enumerate}

\begin{example}[First Order Stationarity]
	If $\mathcal{F}$ is the class of differentiable functions, then we can define 
	\begin{equation}
		\mathcal{L}_f(w) = \|\nabla f(w)\|
	\end{equation}
\end{example}

\begin{definition}[$\epsilon$-Critical Point]
	For differentiable function $f: \mathbb{R}^d \to \mathbb{R}$, $x \in \mathbb{R}^d$ is an \textbf{$\epsilon$-critical point} if 
	\begin{equation}
		\|\nabla f(x)\| \leq \epsilon
	\end{equation}
\end{definition}

\begin{example}[Function Suboptimality]
	For a given function $f$ where $\inf f$ is finite, we can define 
	\begin{equation}
		\mathcal{L}_f (w) = f(w) - \inf f
	\end{equation}
\end{example}

\begin{example}[Distance to the Minimizer]
	Assuming a minimizer exists for $f: S \to \mathbb{R}$, let $\argmin{f}$ be the set of all minimizers. Then, we can define 
	\begin{equation}
		\mathcal{L}_f (w) = d(w, \argmin{f})
	\end{equation}
	For some metric $d$ defined on $S$.  
\end{example}

What we hope to do is find bounds for all three performance criterions for various algorithms. 

To analyze an optimiztion algorithm, we want to see its broad performance over not just one function, but a \textit{class of functions}. For a fixed algorithm, the worst-case iteration complexity is the most popular. 

\begin{definition}[Worst-Case Iteration Complexity]
	Given function class $\mathcal{F}$, suppose an algorithm $\mathcal{A}$ generates a sequences of iterates $w_0 (f), w_1 (f), w_2 (f), \ldots$ for each $f \in \mathcal{F}$. Define the first iteration at which the algorithm achieves accuracy $\epsilon$ as 
	\begin{equation}
		T_\epsilon(A, f) \colon \inf \{ k \geq 0 \colon \mathcal{L}_f (w_k (f)) \leq \epsilon \}
	\end{equation}
	Then, the 
	\begin{enumerate}
	  \item \textbf{worst-case iteration complexity} is the maximum number of iterations required for algorithm $\mathcal{A}$ to achieve accuracy $\epsilon$ over the entire function class $\mathcal{F}$
			\begin{equation}
				K_\epsilon(A, \mathcal{F}, \mathcal{L}) \coloneqq \sup_{f \in \mathcal{F}} T_\epsilon (A, f)
			\end{equation}

		\item \textbf{uniform worst-case iteration complexity} is defined 
			\begin{equation}
				\widetilde{K}_\epsilon(A, \mathcal{F}, \mathcal{L}) \coloneqq \inf \big\{ k \geq 0 \colon \sup_{f \in \mathcal{F}} \mathcal{L}_f (w_k (f)) \leq \epsilon \big\}
			\end{equation}
	\end{enumerate}
\end{definition}


\begin{definition}[Minimax Oracle Error]
	Let $\mathcal{F}$ be a function class, $\mathcal{L}$ a performance criterion, and $\mathcal{A}_k (\mathcal{O})$ denote a specified class of algorithms making at most $k$ calls to oracle $\mathcal{O}$. The \textbf{minimax oracle error} is the smallest worst case error achievable by any admissible algorithm after $k$ oracle queries. 
	\begin{equation}
		\mathcal{E}_k (\mathcal{F}, \mathcal{O}, \mathcal{L}) \coloneq \inf_{A \in \mathcal{A}_k (\mathcal{()})} \sup_{f \in \mathcal{F}} \mathcal{L}_f \big(w_{\mathcal{A}_k} (f) \big)
	\end{equation}
\end{definition}


% 1. Smooth Optimization - gradient descent, L-smooth, \epsilon-critical, Taylor 1/2nd order bound, PL-inequality, linear convergence of GD with PL, Lower bounds on convergence of GD
% 1.1 Convexity - Strong convexity, SC implies PL, convergence of GD on strongly convex functions, ridge regularization to attain strong convexity. 
% 1.2 Momentum/Acceleration - Polyak heavy-ball, nesterov, convergence, 
%
% 2. Nonsmooth Optimization - subgradients, subdifferentials
% 2.1 Convex Conjugacy - Conjugates, Fenchel-Young Inequality, Fenchel-Moreau theorem, Fenchel Duality
% 2.2 Subgradient Descent - convergence theorems
% 2.3 Proximal Methods - ISTA, FISTA, L1 regularization, sparsity, convergence, Moreau Envelope/Decomposition
%
% 3. Second Order Optimization - Newton Raphson Method
% 3.1 Newton's Method - Netwon, Newton Raphson
% 3.2 Quasi-Newton Methods - Secant, BFGS, L-BFGS 
% 3.3 Preconditioned Methods - Preconditioned GD, Condition Numbers, Matrix Preconditioners 
%
% 4. Constrained Optimization - Constrained Optimality 
% 4.1 Lagrangian Duality and KKT - Convex Duality
% 4.2 Augmented Lagrangian Methods - ADMM
% 4.3 Projected Gradient Descent 
% 4.4 Interior Point Methods
%
% 5. Stochastic Optimization - SGD, ERM, unbiasedness, variance of stochastic gradients
% 5.1 Convergence - Smooth, convex, PL, step sizes 
% 5.2 Momentum and Adaptive Methods - Momentum, Adagrad, RMSProp, Adam, Convergence results. 
% 5.3 Variance Reduction 
% 5.4 Proximal Methods - Proximal SGD, Stochastic subgradient descent, Projected SGD
%
% 6. Gradient Free Methods 
% 6.1 Simulated Annealing 
% 6.2 Nelder Mead
